\documentclass{leveldoc}
\DocUses{foundations,slingshot}
\DocFingerprint{f58964bda6e3bdd4}

\begin{document}

\section{What the level asks}\label{sec:asks}

The game's text:
\leveldescription

A positive particle starts inside a large circular coil, at $(\val{xZero}, \val{yZero})$,
moving straight up ($+y$). The coil lies in the plane of the board, centred at
$(\val{xCoil}, \val{yCoil})$ with radius $\val{aCoil}$; its current makes a magnetic field
perpendicular to the board, which bends the particle into a circle. The particle must enter
the detector, the rectangle $\val{xDetLo} \le x \le \val{xDetHi}$,
$\val{yDetLo} \le y \le \val{yDetHi}$, with no acceptance window, no kinetic window and no
gate. The player places up to three charges, of either sign, from the level's magnitudes. A
flight ends when the particle enters the detector, touches a charge or the coil's wire, leaves
the world, or reaches $t_{\max}$ (\xsecref{foundations}{sec:board}).

With nothing placed the particle circles inside the coil, the circle precessing slowly, and
never meets the detector before $t_{\max}$. The reference places one positive charge below the
coil's centre line: its push walks the circling particle across the board, turn by turn, until
it enters the detector. This level opens the magnetic arc, so this document derives the
magnetic force's main facts (it does no work; it bends a charge into a circle of radius
$p/(qB)$) and the field of a circular coil in its own plane, from the Biot--Savart law to a
closed form, and integrates the flight numerically as \levelref{slingshot} does.

\section{Setup}\label{sec:setup}

The level file, field by field (the appendix reproduces it):
\setuptable

In symbols: the particle has charge $q = \val{q}$, mass $m = \val{m}$, radius $a_\mathrm{p} =
0$; the speed of light is $c = \val{c}$. It starts at $\vect{x}_0 = (\val{xZero}, \val{yZero})$
with kinetic energy $T_0 = \val{Tzero}$ along the level file's direction, the computed
$(\cos 90^\circ, \sin 90^\circ)$, whose first component $\val{dirX}$ is a rounding error. The
coil is a circle of radius $a = \val{aCoil}$ about $\vect{r}_\mathrm{c} = (\val{xCoil},
\val{yCoil})$, carrying the current given by its strength $\kappa = \mu_0 I/4\pi =
\val{kappa}$ (field units, \physref{2.2}); its wire has radius $\val{wire}$. The reference
charge is $Q = \val{QOne}$ at $(\val{xOne}, \val{yOne})$, radius $R = \val{R}$, coupling
$\kappa_Q = qQ = \val{kOne}$ (repulsive). The world is $-1 \le x \le \val{xWorld}$,
$-1 \le y \le \val{yWorld}$.

\begin{figure}[htbp]
\centering
\begin{tikzpicture}
\begin{axis}[axis equal image, width=0.92\linewidth, xmin=-1.6, xmax=31.6, ymin=-1.6, ymax=21.6,
  xtick={0,5,...,30}, ytick={0,5,...,20}, tick label style={font=\small},
  xlabel={$x$}, ylabel={$y$}, label style={font=\small}]
  \draw[black!8] (0,0) grid (30,20);
  \draw[black!45, dashed] (-1,-1) rectangle (\rawval{xWorld},\rawval{yWorld});
  \draw[black!70, thick] (\rawval{xCoil},\rawval{yCoil}) circle[radius=\rawval{aCoil}];
  \fill[green!50!black, fill opacity=0.18] (\rawval{xDetLo},\rawval{yDetLo}) rectangle (\rawval{xDetHi},\rawval{yDetHi});
  \draw[green!40!black] (\rawval{xDetLo},\rawval{yDetLo}) rectangle (\rawval{xDetHi},\rawval{yDetHi});
  \addplot[orange!70!black, very thin] table {pathbare.dat};
  \addplot[blue!70!black] table {path.dat};
  \fill[red!75!black] (\rawval{xOne},\rawval{yOne}) circle[radius=\rawval{R}];
  \fill (\rawval{xZero},\rawval{yZero}) circle[radius=0.15];
\end{axis}
\end{tikzpicture}
\caption{The coil (black circle) and the two flights. Blue: the reference, whose circles walk
across the board under the push of the charge below (red) until they enter the detector (green).
Orange: the empty board, circling and precessing inside the coil until $t_{\max}$.}
\label{fig:board}
\end{figure}

\section{Principles and model}\label{sec:model}

\begin{principles}
  \item The Lorentz force $\vect{F} = q(\vect{E} + \vect{v}\times\vect{B})$ in the
    relativistic equation of motion \xeqref{foundations}{eq:eom}, with the launch
    \xeqref{foundations}{eq:launch}.
  \item The field of a point charge \xeqref{foundations}{eq:coulomb}; the conservation of energy
    \xeqref{foundations}{eq:energy}, which a static magnetic field leaves unchanged.
  \item The Biot--Savart law for a closed current, $\vect{B}(\vect{x}) =
    \kappa\oint \dd\vect{l}\times(\vect{x} - \vect{x}')/\abs{\vect{x} - \vect{x}'}^3$ with
    $\kappa = \mu_0 I/4\pi$ (\Jackson{§5.2}), and the game's units, $\mu_0/4\pi = 1/c^2$
    (\physref{2.2}).
  \item Complete elliptic integrals $K(m)$, $E(m)$, their derivatives and the
    arithmetic--geometric mean (DLMF §19.4, §19.8).
  \item The board's rules (\xsecref{foundations}{sec:board}); the classical Runge--Kutta method,
    cubic Hermite interpolation and the error estimate by doubling the step
    (\xsecref{foundations}{sec:numerics}), as in \xsecref{slingshot}{sec:method}.
\end{principles}

In the plane of the coil its field is perpendicular to the plane, $\vect{B} = B_z\uvec{z}$
(\S\ref{sec:loop}), so with $\vect{v} = (v_x, v_y, 0)$ the magnetic force
$q\vect{v}\times\vect{B} = qB_z(v_y, -v_x, 0)$ stays in the plane, and the motion with it. The
state $(\vect{x}, \vect{p})$ obeys
\begin{equation}\label{eq:eom-b}
  \deriv{\vect{x}}{t} = \vect{v} = \frac{\vect{p}}{\gamma m}, \qquad
  \deriv{\vect{p}}{t} = \kappa_Q\frac{\vect{x} - \vect{r}_Q}{\abs{\vect{x} - \vect{r}_Q}^3}
    + qB_z(\vect{x})\,(v_y, -v_x), \qquad
  \gamma = \sqrt{1 + \frac{\abs{\vect{p}}^2}{m^2c^2}},
\end{equation}
from $\vect{x}_0$ and $\vect{p}_0 = \abs{\vect{p}_0}(\val{dirX}, 1)$ until the first event.

\begin{assumption}{The coil and the charge are held in place}\label{as:fixed}
Neither recoils, and the coil's current is fixed (rules of the game).
\end{assumption}

\begin{assumption}{No radiation}\label{as:radiation}
The radiated energy \xeqref{foundations}{eq:larmor} is not taken from the particle. Along the
reference flight it is $\val{Erad}$, a fraction $\val{EradFrac}$ of $T_0$.
\end{assumption}

\begin{assumption}{A thin coil in the plane}\label{as:wire}
The current runs along the circle of radius $a$ itself; the wire's radius $\val{wire}$ only
makes it an obstacle. The game does the same (\physref{2.2}).
\end{assumption}

\section{Derivation}\label{sec:derivation}

\subsection{The magnetic force does no work}\label{sec:nowork}

With $\vect{E} = 0$, $\dd\abs{\vect{p}}^2/\dd t = 2\vect{p}\cdot q(\vect{v}\times\vect{B}) = 0$,
since $\vect{p} \parallel \vect{v}$. So $\abs{\vect{p}}$, $\gamma$ and the speed are constant:
the field turns the velocity and never changes its size. With the charge's field as well, the
energy \xeqref{foundations}{eq:energy},
\begin{equation}\label{eq:energy-b}
  W = \gamma mc^2 + \frac{\kappa_Q}{\abs{\vect{x} - \vect{r}_Q}},
\end{equation}
is conserved: $\vect{B}$ adds nothing to it. At the launch $V(\vect{x}_0) = \val{VZero}$ and
$W = \val{W}$.

\subsection{Circles}\label{sec:circle}

In a uniform field $B_z$ and with $\vect{E} = 0$, $\abs{\vect{p}}$ is constant and
\eqref{eq:eom-b} turns $\vect{p}$ at the constant rate
\begin{equation}\label{eq:cyclotron}
  \omega_\mathrm{c} = \frac{\abs{q B_z}}{\gamma m}, \qquad
  T_\mathrm{c} = \frac{2\pi}{\omega_\mathrm{c}} = \frac{2\pi\gamma m}{\abs{qB_z}} ,
\end{equation}
($\dot p_x = \omega_\mathrm{c}\,p_y$, $\dot p_y = -\omega_\mathrm{c}\,p_x$ for $qB_z > 0$: a
clockwise rotation). The position follows a circle of radius
\begin{equation}\label{eq:larmor-radius}
  r_\mathrm{L} = \frac{v}{\omega_\mathrm{c}} = \frac{\abs{\vect{p}}}{\abs{q B_z}} .
\end{equation}
The relativistic momentum enters through $\gamma$ in the period; the radius depends on the
momentum alone.

\subsection{The field of a circular coil in its plane}\label{sec:loop}

Put the coil's centre at the origin and the field point on the $x$ axis at the distance $\rho$.
The wire element at the angle $\phi$ is $\dd\vect{l} = a(-\sin\phi, \cos\phi, 0)\,\dd\phi$ at
$\vect{x}' = a(\cos\phi, \sin\phi, 0)$, so $\vect{x} - \vect{x}' = (\rho - a\cos\phi,
-a\sin\phi, 0)$ and $(\dd\vect{l}\times(\vect{x} - \vect{x}'))_z = a(a - \rho\cos\phi)\,\dd\phi$;
the other components vanish in the plane. The Biot--Savart law gives
\begin{equation}\label{eq:biot-savart}
  B_z(\rho) = \kappa\int_0^{2\pi} \frac{a(a - \rho\cos\phi)}{(a^2 + \rho^2 - 2a\rho\cos\phi)^{3/2}}\,\dd\phi ,
\end{equation}
and at the centre $B_z(0) = 2\pi\kappa/a = \val{Bcentre}$. To bring \eqref{eq:biot-savart}
to a closed form, substitute $\phi = \pi - 2\theta$: then
$a^2 + \rho^2 - 2a\rho\cos\phi = \beta^2(1 - m\sin^2\theta)$ and
$a - \rho\cos\phi = a + \rho - 2\rho\sin^2\theta$, with
\[
  \beta = a + \rho, \qquad \alpha^2 = (a - \rho)^2, \qquad m = \frac{4a\rho}{\beta^2},
  \qquad 1 - m = \frac{\alpha^2}{\beta^2}.
\]
The integrand is even in $\theta$, and $\phi \in [0, 2\pi]$ covers $\theta \in [-\pi/2, \pi/2]$
with $\dd\phi = -2\,\dd\theta$, so
\[
  B_z = \frac{4\kappa a}{\beta^3}\int_0^{\pi/2}
    \frac{a + \rho - 2\rho\sin^2\theta}{(1 - m\sin^2\theta)^{3/2}}\,\dd\theta .
\]
Two integrals reduce to $K(m) = \int_0^{\pi/2}(1 - m\sin^2\theta)^{-1/2}\dd\theta$ and
$E(m) = \int_0^{\pi/2}(1 - m\sin^2\theta)^{1/2}\dd\theta$. Differentiating $K$ under the
integral, $2\,\dd K/\dd m = \int_0^{\pi/2}\sin^2\theta\,(1 - m\sin^2\theta)^{-3/2}\dd\theta$, and
$\dd K/\dd m = (E - (1 - m)K)/(2m(1 - m))$ (DLMF §19.4), so
\[
  \int_0^{\pi/2}\frac{\sin^2\theta\,\dd\theta}{(1 - m\sin^2\theta)^{3/2}}
    = \frac{1}{m}\Big(\frac{E}{1 - m} - K\Big), \qquad
  \int_0^{\pi/2}\frac{\dd\theta}{(1 - m\sin^2\theta)^{3/2}} = \frac{E}{1 - m},
\]
the second from the first by writing $1 = (1 - m\sin^2\theta) + m\sin^2\theta$ in the numerator.
Inserting both, with $2\rho/m = \beta^2/(2a)$ and $1/(1 - m) = \beta^2/\alpha^2$,
\begin{align*}
  B_z &= \frac{4\kappa a}{\beta}\Big[\frac{(a + \rho)E}{\alpha^2}
    - \frac{1}{2a}\Big(\frac{\beta^2 E}{\alpha^2} - K\Big)\Big] \\
    &= \frac{2\kappa}{\alpha^2\beta}\big[(2a^2 + 2a\rho - \beta^2)E + \alpha^2 K\big],
\end{align*}
and since $2a^2 + 2a\rho - \beta^2 = a^2 - \rho^2$,
\begin{equation}\label{eq:loop-field}
  B_z(\rho) = \frac{2\kappa}{\alpha^2\beta}\Big[(a^2 - \rho^2)\,E(m) + \alpha^2 K(m)\Big],
\end{equation}
the in-plane case of the closed form the game uses (\physref{2.2}). As $\rho \to 0$,
$K, E \to \pi/2$ and \eqref{eq:loop-field} returns $2\pi\kappa/a$. $K$ and $E$ are computed by
the arithmetic--geometric mean (DLMF §19.8): from $a_0 = 1$, $b_0 = \sqrt{1 - m}$,
$c_0^2 = m$, iterate $a_{n+1} = (a_n + b_n)/2$, $b_{n+1} = \sqrt{a_nb_n}$,
$c_{n+1} = (a_n - b_n)/2$ to convergence (quadratic), then
\begin{equation}\label{eq:agm}
  K = \frac{\pi}{2a_\infty}, \qquad E = K\Big(1 - \sum_{n\ge 0} 2^{n-1}c_n^2\Big).
\end{equation}
Independently, \eqref{eq:biot-savart} is the integral of a smooth periodic function, for which
the trapezoidal rule converges faster than any power of the step; with $\val{nTrap}$ nodes
the two agree to a relative $\val{fieldAgree}$:
\begin{center}\small\val{fieldtab}\end{center}
The field grows from the centre towards the wire and diverges at it ($\alpha \to 0$).

\subsection{The launch}\label{sec:launch}

$\gamma_0 = 1 + T_0/(mc^2) = \val{gammaZero}$, $\abs{\vect{p}_0} = \val{pZero}$,
$v_0 = \val{vZero}$. The launch is $\rho_0 = \val{rhoZero}$ from the coil's centre, where
$B_z = \val{Bzero}$. By \eqref{eq:larmor-radius} and \eqref{eq:cyclotron} the particle would
circle there with $r_\mathrm{L} = \val{rL}$ and $T_\mathrm{c} = \val{Tc}$
($\omega_\mathrm{c} = \val{omegaC}$), \val{sense}, about $(\val{xGuide}, \val{yGuide})$. That
circle reaches the far side of the coil's centre, where the field is about half as strong
(the table): the orbit does not close, and the motion needs \eqref{eq:eom-b} itself.

\subsection{The method}\label{sec:method}

The state is advanced by \eqref{eq:eom-b} with the classical Runge--Kutta method
\xeqref{foundations}{eq:rk4} at the fixed step $h = \val{H}$, $B_z$ from \eqref{eq:loop-field}
at every evaluation. After each step the cubic Hermite interpolant of the positions
(\xeqref{slingshot}{eq:hermite}) is sampled at 64 points; the first sample inside the detector,
the charge's disc of radius $R$, within $\val{wire}$ of the coil's circle, or outside the world
is refined by bisection, and the state there comes from one more step. A flight still going at
$t_{\max}$ stops there, its last step shortened to end on $t_{\max}$. The error is of order
$h^4$, estimated by doubling $h$ (\xsecref{foundations}{sec:numerics}); the energy
\eqref{eq:energy-b} checks every step.

\section{The reference solution by hand}\label{sec:hand}

\begin{procedure}
  \item\label{st:launch} \emph{The launch.} $\gamma_0 = \val{gammaZero}$,
    $\abs{\vect{p}_0} = \val{pZero}$, $\vect{p}_0 = \abs{\vect{p}_0}(\val{dirX}, 1)$;
    $W = \val{W}$.
  \item\label{st:field} \emph{The field.} At each evaluation, $\rho = \abs{\vect{x} -
    \vect{r}_\mathrm{c}}$, $m = 4a\rho/(a + \rho)^2$, $K$ and $E$ by \eqref{eq:agm}, $B_z$ by
    \eqref{eq:loop-field}.
  \item\label{st:march} \emph{The march.} Advance \eqref{eq:eom-b} by
    \xeqref{foundations}{eq:rk4}. The state every five time units (step $\val{HTable}$; the
    last column the energy's drift):
    \begin{center}\small\val{states}\end{center}
  \item\label{st:entry} \emph{The entry.} With $h = \val{H}$, after $\val{steps}$ steps, the
    particle enters the detector through its \val{faceE} face at $t_\mathrm{e} = \val{tE}$, at
    the point $(\val{xE}, \val{yE})$ and with the momentum
    $\vect{p}_\mathrm{e} = (\val{pxE}, \val{pyE})$. Doubling the step estimates the error of
    that end state as $\val{err}$.
  \item\label{st:checks} \emph{The checks.} $W$ drifts by at most $\val{wDev}$. From $W$ alone,
    $\abs{\vect{p}_\mathrm{e}} = \val{pEnergy}$, which the march matches to $\val{pDiff}$.
  \item\label{st:rad} \emph{The radiation.} Simpson's rule on the steps gives $\val{Erad}$;
    against the step $2h$ its error is estimated as a relative $\val{EradErr}$.
\end{procedure}

\begin{result}[The reference flight]
The particle enters the detector through its \val{faceE} face at $(\val{xE}, \val{yE})$ at
$t_\mathrm{e} = \val{tE}$, with momentum $(\val{pxE}, \val{pyE})$. The level is solved.
\end{result}

\section{Comparison with the game}\label{sec:comparison}

\comparisontable

Both flights agree to a few $10^{-10}$ after many turns: the reference's time by
$\val{cmpTDiff}$, the empty board's position at $t_{\max}$ by $\val{cmpBareXDiff}$, the size of
the game's verified accuracy (Foundations §6) accumulated over $\val{bareTurns}$ turns. Our own
step estimates are smaller ($\val{err}$ and $\val{bareErr}$). Each flight's tolerance is the
larger of $10^{-9}$, ten times our estimate, and the game's own gap between its preview and its
verification ($\val{gapRef}$ and $\val{gapBare}$), fixed before comparing. The radiated energies
agree to $\val{cmpRadDiff}$ (relative).

\section{Other solutions}\label{sec:others}

The empty board, the second placement in \code{placements.json}, does not solve the level: the
particle circles inside the coil for $\val{bareTurns}$ cyclotron periods and is still flying at
$t_{\max}$, at $(\val{bareX}, \val{bareY})$; its energy drifts by $\val{bareWdev}$. Its circle
precesses because the field varies across it (\S\ref{sec:launch}). A charge works by pushing the
circle's centre: the reference's repulsion from below moves it step by step until the particle
enters the detector, after many turns. Charges of either sign and other positions do the same
in other ways; a map of them would take the game's verdicts, which this document does not
reproduce.

\section{What the game computes}\label{sec:engine}

The game evaluates the coil's field by the closed form of Simpson et al.\ with $K$ and $E$ from
the arithmetic--geometric mean (\physref{2.2}), of which \eqref{eq:loop-field} is the case
$z = 0$, and integrates \eqref{eq:eom-b} with the adaptive eighth-order Runge--Kutta method of
Dormand and Prince (\physref{5}), with its event location (\physref{6}) and verification
(\physref{7}). In the plane its in-plane $\vect{B}$ components are exactly zero, so the
particle stays in the plane, as here. The differences in the comparison are the two
integrations' errors. Radiation is integrated, not applied (\physref{8}).

\section{Further reading}\label{sec:reading}

The Biot--Savart law and the circular loop's field off its axis, through its vector potential
and elliptic integrals: \Jackson{§5.2, §5.5}. Motion in a magnetic field, the cyclotron
frequency and drifts in a slowly varying field: \Jackson{§12.2, §12.4}. Elliptic integrals and the
arithmetic--geometric mean: DLMF Ch.~19. The numerical method: \levelref{slingshot}.

\levelfileappendix

\end{document}
